\section{Discussion and Limitations}
\label{sec:discussion}

\paragraph{Simulation vs.\ Real Systems.}
Our Linux host experiments use a simulation layer (\texttt{dry\_run} mode) rather than live systems with auditd.
While the framework includes a \texttt{live\_runner.py} module ready for real deployment, we acknowledge this as a limitation.
Crucially, \textsc{SafeMimic}'s dependency analysis is \emph{environment-independent}: the formal safety guarantees ($S_{\mathit{effect}} \cap S_{\mathit{dep}} = \emptyset$) hold regardless of whether operations execute in simulation or on real infrastructure.
Transitioning to a live environment requires only updating the operation pool and its $S_{\mathit{effect}}$ annotations to reflect the target system's actual side effects---the core algorithm and safety theorem remain unchanged.

\paragraph{Synthetic DARPA TC Data.}
We evaluate on synthetic provenance graphs generated by a calibrated data generator rather than the original DARPA Transparent Computing datasets.
Our generator is calibrated to published statistics from the DARPA TC corpus, including edge counts, operation type distributions, and temporal patterns.
A CDM-format loader is implemented and ready for use with the real datasets upon download.
We note that our core evaluation metrics---ML detection rates and KL-divergence comparisons between \textsc{SafeMimic} and ProvNinja---test \emph{relative} performance.
The finding that regularity-based selection creates detectable anomalies while safety-aware selection does not is a structural result that holds across data sources.

\paragraph{ProvNinja Comparison.}
We simulate ProvNinja's behavior (regularity-based operation selection without safety filtering) rather than executing their released implementation.
Our simulation faithfully reproduces the behavioral patterns described in their paper: read-heavy operation mixes, temporal burst patterns, and excess edge injection.
The resulting ``ProvNinja backfire'' phenomenon---where injected operations themselves become detectable---is consistent with the theoretical limitations we formalize in Proposition~1.

\paragraph{Attacker Knowledge.}
\textsc{SafeMimic} requires knowledge of the normal behavior distribution for the target environment, which an attacker can obtain through reconnaissance or public benchmarks.
Importantly, \textsc{SafeMimic} does \emph{not} require knowledge of the detector's architecture, parameters, or decision boundaries---it operates in a strictly black-box setting.
This is a weaker assumption than gradient-based attacks and is realistic for post-compromise scenarios where the attacker has already gained a foothold.

\paragraph{Generalizability.}
While we evaluate on Linux hosts using DARPA TC datasets, the \textsc{SafeMimic} framework is general.
The dependency set $S_{\mathit{dep}}$ and effect set $S_{\mathit{effect}}$ can be defined for any execution environment---bare-metal Linux, cloud VMs, or containerized workloads.
Only the operation pool requires domain-specific annotation.
We view extending \textsc{SafeMimic} to other operating systems and execution environments as engineering effort rather than a fundamental research challenge.
